% 复数乘法的伸缩旋转与三次单位根。解析数据：z=1+2i，w=3-i，zw=5+5i；
% arg z=atan2(2,1)≈63.435°，arg w=atan2(-1,3)≈-18.435°，arg(zw)=45°。
\begin{tikzpicture}[line width=.6pt,font=\fontsize{8}{10}\selectfont,>=Stealth]
% ---------- (a) 乘法 ----------
\begin{scope}[x=7.2mm,y=7.2mm]
 \draw[FigureGrid,line width=.4pt] (-1,-1.5) grid (5.6,5.6);
 \draw[->,FigureStroke,line width=.75pt] (-1,0)--(5.9,0) node[right,FigureInk]{$\operatorname{Re}$};
 \draw[->,FigureStroke,line width=.75pt] (0,-1.5)--(0,5.9) node[above,FigureInk]{$\operatorname{Im}$};
 % 三个辐角
 \draw[FigureBlue,line width=.6pt] (1.1,0) arc[start angle=0,end angle=63.435,radius=1.1];
 \node[FigureBlue] at (1.1,1.4) {$\theta_1$};
 \draw[FigureGreen,line width=.6pt] (1.7,0) arc[start angle=0,end angle=-18.435,radius=1.7];
 \node[FigureGreen] at (2.05,-0.75) {$\theta_2$};
 \draw[FigurePurple,line width=.6pt] (2.4,0) arc[start angle=0,end angle=45,radius=2.4];
 \node[FigurePurple,right] at (2.35,0.85) {$\theta_1+\theta_2$};
 % 三个点与向径
 \draw[FigureBlue,line width=1.1pt] (0,0)--(1,2);
 \fill[FigureBlue] (1,2) circle[radius=.9mm];
 \node[FigureBlue,left] at (0.95,2.25) {$z=1+2i$};
 \draw[FigureGreen,line width=1.1pt] (0,0)--(3,-1);
 \fill[FigureGreen] (3,-1) circle[radius=.9mm];
 \node[FigureGreen,right] at (3.05,-1.05) {$w=3-i$};
 \draw[FigurePurple,line width=1.1pt] (0,0)--(5,5);
 \fill[FigurePurple] (5,5) circle[radius=.9mm];
 \node[FigurePurple,left] at (4.9,5.25) {$zw=5+5i$};
 \node[FigureInk] at (2.3,-2.3) {(a) 模相乘、辐角相加};
\end{scope}
% ---------- (b) 三次单位根 ----------
\begin{scope}[xshift=75mm,yshift=17mm,x=17mm,y=17mm]
 \draw[->,FigureStroke,line width=.75pt] (-1.45,0)--(1.5,0) node[right,FigureInk]{$\operatorname{Re}$};
 \draw[->,FigureStroke,line width=.75pt] (0,-1.35)--(0,1.45) node[above,FigureInk]{$\operatorname{Im}$};
 \draw[FigureMuted,dashed,line width=.6pt] (0,0) circle[radius=1];
 \fill[FigureBlueFill] (1,0)--(-0.5,0.8660)--(-0.5,-0.8660)--cycle;
 \draw[FigureBlue,line width=.75pt] (1,0)--(-0.5,0.8660)--(-0.5,-0.8660)--cycle;
 \draw[FigureGreen,line width=.6pt] (0.28,0) arc[start angle=0,end angle=120,radius=0.28];
 \node[FigureGreen] at (0.36,0.2) {$\tfrac{2\pi}{3}$};
 \foreach \a in {0,120,240}{\fill[FigureBlue] (\a:1) circle[radius=.9mm];}
 \node[FigureInk,above right] at (1,0) {$1$};
 \node[FigureInk,above left] at (-0.5,0.8660) {$-\tfrac12+\tfrac{\sqrt3}{2}i$};
 \node[FigureInk,below left] at (-0.5,-0.8660) {$-\tfrac12-\tfrac{\sqrt3}{2}i$};
 \node[FigureInk] at (0,-1.75) {(b) $z^3=1$的三个根};
\end{scope}
\end{tikzpicture}
